% from ega4-11-fr.tex
\oldpage[IV-3]{142}

il faut et il suffit que $P\otimes_C(C/\mathfrak N)$ soit un
$(C/\mathfrak N)$-module libre. Notons d'autre part que si $C$ est un anneau
noethérien réduit, il existe $g\in C$ tel que $C_g$ soit un anneau intègre.
Utilisons maintenant le lemme (6.9.2)~: on peut définir par récurrence une suite
$(f_i)_{i\geq1}$ d'éléments de $A$ de la façon suivante~:
\begin{enumerate}
\item[$1^\circ$] $f_1$ est tel que $A_1=(A_{\mathrm{red}})_{f_1}$ soit
intègre et $M\otimes_A A_1$ un $A_1$-module libre~;
\item[$2^\circ$] si l'idéal $\mathfrak J_i$ engendré par $f_1,\ldots,f_i$
est $A$, $f_{i+1}=f_i$~;
\item[$3^\circ$] si au contraire $\mathfrak J_i\neq A$, $f_{i+1}$ est un
élément n'appartenant pas à $\mathfrak J_i$ tel que
$A_{i+1}=((A/\mathfrak J_i)_{\mathrm{red}})_{f_{i+1}}$ soit intègre et
$M\otimes_A A_{i+1}$ un $A_{i+1}$-module libre.
\end{enumerate}
Comme $A$ est noethérien, la suite croissante $(\mathfrak J_i)$ est
stationnaire, donc il existe $n$ tel que $f_1,\ldots,f_n$ engendrent l'idéal
$A$, et les $f_i$ répondent à la question puisque
\[
((A/\mathfrak J_i)_{\mathrm{red}})_{f_{i+1}}
=((A/\mathfrak J_i)_{f_{i+1}})_{\mathrm{red}}.
\]
\end{proof}

\begin{proposition}[11.3.16]
\label{IV.11.3.16-fr}
Soient $f:X\to Y$ un morphisme fidèlement plat de présentation finie,
$g:Y\to Z$ un morphisme tel que $g\circ f:X\to Z$ soit un morphisme de type
fini (resp. de présentation finie). Alors $g$ est un morphisme de type fini
(resp. de présentation finie).
\end{proposition}

\begin{proof}
Comme $f$ est surjectif et $g\circ f$ quasi-compact, $g$ est quasi-compact
(1.1.3). Montrons d'abord que si $g\circ f$ est de présentation finie, $g$ est
quasi-séparé. En effet, considérons un ouvert affine $W\subset Z$, et soit
$(V_i)$ un recouvrement affine fini de $g^{-1}(W)$~; il s'agit de voir que les
$V_i\cap V_j$ sont quasi-compacts (1.2.6 et 1.2.7). Pour chaque $i$, soit
$(U_{ih})$ un recouvrement ouvert affine fini de $f^{-1}(V_i)$~; comme $f$ est
de présentation finie, les $U_{ih}\cap U_{jk}$ sont tous quasi-compacts~; or,
puisque $f$ est surjectif, $V_i\cap V_j$ est réunion des images
$f(U_{ih}\cap U_{jk})$ pour $h$, $k$ variables, donc est quasi-compact puisque
$f$ est continue.

La question est donc locale sur $Y$ et on peut supposer que
$Y=\operatorname{Spec}(B)$ et $Z=\operatorname{Spec}(A)$ sont affines, $X$
étant réunion finie d'ouverts affines $X_i=\operatorname{Spec}(C_i)$, où les
$C_i$ sont des $B$-algèbres de type fini (resp. de présentation finie). Si
$X'$ est le préschéma somme des $X_i$, $p:X'\to X$ le morphisme qui coïncide
avec l'injection canonique dans chaque $X_i$, $p$ est un morphisme fidèlement
plat de présentation finie (1.6.5), donc $g\circ f\circ p$ est un morphisme de
type fini (resp. de présentation finie) et $f\circ p$ un morphisme fidèlement
plat de présentation finie. On peut par suite supposer aussi que
$X=\operatorname{Spec}(C)$ est affine, et on est donc ramené à prouver le
corollaire suivant.
\end{proof}

\begin{corollary}[11.3.17]
\label{IV.11.3.17-fr}
Soient $A$ un anneau, $B$ une $A$-algèbre, $C$ une $B$-algèbre de présentation
finie et qui soit un $B$-module fidèlement plat. Alors, si $C$ est une
$A$-algèbre de type fini (resp. de présentation finie), $B$ est une
$A$-algèbre de type fini (resp. de présentation finie).
\end{corollary}

\begin{proof}
Supposons d'abord que $g\circ f$ soit de type fini. Soit $(B_\lambda)$ la
famille filtrante croissante des sous-$A$-algèbres de type fini de $B$~; en
vertu de (8.8.2), il existe un indice $\alpha$ tel que
$C=C_\alpha\otimes_{B_\alpha}B$, où $C_\alpha$ est une
$B_\alpha$-algèbre de présentation finie~; en outre (8.10.5 et 11.2.6) on peut
supposer que $C_\alpha$ est un $B_\alpha$-module fidèlement plat. Pour
$\lambda\geq\alpha$, on a donc
$C=C_\lambda\otimes_{B_\lambda}B$, où $C_\lambda$ est un
$B_\lambda$-module fidèlement plat~; comme l'application $B_\lambda\to B$ est
injective, on en déduit qu'il en est de même de $C_\lambda\to C$. En outre,
comme $C=\varinjlim C_\lambda$, tout élément de $C$ appartient à un
$C_\lambda$, et par suite il existe $\lambda$ tel que l'application
$C_\lambda\to C$

